%%%%%%%%%%%%%%%%%%%%%%%%%%%%%%%%%%%%%%%%%
% 模板資訊：
% 模板名稱：Ernie's Beamer
% 版本：1.0 (2023.03.24)
% 修改者：莊程翔 Ernie Cheng-Xiang Zhuang
% 編譯器：XeLaTeX
%
% 原始模板的資訊：
% 模板名稱：Beamer Presentation
% 作者：Vel (vel@latextemplates.com)
% 編譯器：XeLaTeX
% 授權：CC BY-NC-SA 4.0 (https://creativecommons.org/licenses/by-nc-sa/4.0/)
% 下載連結：https://www.LaTeXTemplates.com
%
% 製作本模板之目的：
% 由於，LaTeX 對初學者來說不太友善，因此我針對 Vel 製作的模板做了大幅度的修改及附上清楚明瞭的註解，同時添加了許多常用且好用的封包，並增添了許多常用的客製化設定指令，以便 LaTeX 初學者能夠輕鬆地完成專業的學術簡報。
% 如果您有任何問題，可以透過以下兩種方式聯繫我：
% 1. 網站：https://www.ernie-zhuang.com/contact
% 2. Email：erniezhuang1127@gmail.com
%%%%%%%%%%%%%%%%%%%%%%%%%%%%%%%%%%%%%%%%%

%----------------------------------------------------------------------------------------
%	封包與文檔配置
%----------------------------------------------------------------------------------------

% 自訂字體的封包
\usepackage{fontspec} 

%% 設定中文字體
\usepackage{xeCJK}
\setCJKmainfont[AutoFakeBold=3.5 , AutoFakeSlant=0.2, Path = fonts/]{cwTeX_圓體.ttf}
\XeTeXlinebreaklocale "zh"
\XeTeXlinebreakskip = 0pt plus 1pt

%% 設定英文字體
\usefonttheme{serif} % 使用 serif 字體

% 自訂字體顏色的封包
\usepackage{color} 

%% 自訂顏色
\definecolor{NTHU_Purple}{RGB}{126,35,138}
\definecolor{Default_Blue}{RGB}{52,51,171}

% 數學工具及符號
\usepackage{mathtools, amsmath, amsfonts, amsthm, latexsym} 

% 分別將數學符號間的間隔加大及加粗
\usepackage{newtxtext,newtxmath}

% 圖表自動編號的封包
\usepackage[justification=centering]{caption} 

%% 設定自動編號
\setbeamertemplate{caption}[numbered]

%% 設定圖表編號及標籤的字體大小及字形
\captionsetup[figure]{font=small, labelfont=md}
\captionsetup[table]{font=small, labelfont=md}

% 導入圖形與表格的封包
\usepackage{graphicx}  % \scalebox{} 可用於將過大的表格縮小
\usepackage{booktabs}

% 排列多個子圖形的封包
\usepackage{subfigure} 

% 允許表格的一格能多列呈現的封包
\usepackage{multirow} 

% 可指定表格排版的封包
\usepackage{array}

% 序列標號
\usepackage{enumerate} 

% 繪圖封包 (用於添加浮水印)
\usepackage{tikz}

% 引注參考資料
\usepackage{natbib}

% 註釋掉大部分的封包
\usepackage{comment}

% 設定中文的標籤
\renewcommand{\figurename}{圖} 
\renewcommand{\tablename}{表} 

%----------------------------------------------------------------------------------------
%	排版形式 (擇一，不選等同選擇默認的排版形式)
%----------------------------------------------------------------------------------------

\mode<presentation>{
%\usetheme{default}
%\usetheme{AnnArbor}
%\usetheme{Antibes}
%\usetheme{Bergen}
%\usetheme{Berkeley}
%\usetheme{Berlin}
%\usetheme{Boadilla}
%\usetheme{CambridgeUS}
%\usetheme{Copenhagen}
%\usetheme{Darmstadt}
%\usetheme{Dresden}
%\usetheme{Frankfurt}
%\usetheme{Goettingen}
%\usetheme{Hannover}
%\usetheme{Ilmenau}
%\usetheme{JuanLesPins}
%\usetheme{Luebeck}
\usetheme{Madrid}
%\usetheme{Malmoe}
%\usetheme{Marburg}
%\usetheme{Montpellier}
%\usetheme{PaloAlto}
%\usetheme{Pittsburgh}
%\usetheme{Rochester}
%\usetheme{Singapore}
%\usetheme{Szeged}
%\usetheme{Warsaw}

%----------------------------------------------------------------------------------------
%	外框形式 (擇一，不選等同選擇默認的外框形式)
%----------------------------------------------------------------------------------------

%\useoutertheme{default}
%\useoutertheme{infolines}
%\useoutertheme{miniframes}
\useoutertheme{smoothbars}
%\useoutertheme{sidebar}
%\useoutertheme{split}
%\useoutertheme{shadow}
%\useoutertheme{tree}
%\useoutertheme{smoothtree}

%----------------------------------------------------------------------------------------
%	外框的自訂義調整 
%----------------------------------------------------------------------------------------

% 外框上緣的字 (fg) 為黑色，背景 (bg) 為白色。
%\setbeamercolor{section in head/foot}{fg=black, bg=white} 

% 外框上緣顯示的章節(section)頁數標籤是否關閉
%\setbeamertemplate{mini frames}{}  

% 調整外框形式的字體大小
\setbeamerfont{headline}{size=\scriptsize}
\setbeamerfont{footline}{size=\scriptsize}

% 取消右下方的跳轉工具列
\setbeamertemplate{navigation symbols}{} 

%% 自定義1：外框下緣僅出現名字及頁碼
%\setbeamertemplate{footline} 
%{\leavevmode%
%\hbox{%
%\begin{beamercolorbox}[wd=0.5\paperwidth,ht=3ex,dp=1ex,leftskip=3ex]%
%{author in head/foot}%
%{\scriptsize\textbf{\insertshortauthor}}%
%\end{beamercolorbox}%
%\begin{beamercolorbox}[wd=0.5\paperwidth,ht=3ex,dp=1ex,right]%
%{author in head/foot}%
%\scriptsize \textbf{{\insertframenumber{} / \inserttotalframenumber\hspace*{2ex}}} %頁碼控制選項
%\end{beamercolorbox}%
%}}

%% 自定義2：清除外框下緣但僅出頁碼
%\setbeamertemplate{footline}[page number] 

%% 自定義3：清除外框下緣
%\setbeamertemplate{footline}[] 

%----------------------------------------------------------------------------------------
%	顏色主題 (擇一，不選等同選擇默認的顏色主題)
%----------------------------------------------------------------------------------------

\usecolortheme{default}
%\usecolortheme{albatross}
%\usecolortheme{beaver}
%\usecolortheme{beetle}
%\usecolortheme{crane}
%\usecolortheme{dolphin}
%\usecolortheme{dove}
%\usecolortheme{fly}
%\usecolortheme{lily}
%\usecolortheme{orchid}
%\usecolortheme{rose}
%\usecolortheme{seagull}
%\usecolortheme{seahorse}
%\usecolortheme{whale}
%\usecolortheme{wolverine}

%----------------------------------------------------------------------------------------
%	顏色主題的自訂義調整 
%----------------------------------------------------------------------------------------

% 全文的主題色 (可以特別針對報告對象或機構的代表色調整！)
\setbeamercolor{structure}{fg=NTHU_Purple} 

% 封面頁中標題區塊的底色及字體顏色
%\setbeamercolor{title}{bg=green, fg=black} 

% 各頁標題區塊的底色及字體顏色
%\setbeamercolor{frametitle}{bg=green,fg=black} 

% 全文的內文顏色
%\setbeamercolor{normal text}{fg=orange}

% 數學區塊的標題顏色 
%\setbeamercolor{block title}{bg=blue,fg=yellow} 

% 數學區塊的內文顏色 
%\setbeamercolor{block body}{bg=green,fg=red} 

% 警示文字的顏色
\setbeamercolor{alerted text}{fg=red} 

%----------------------------------------------------------------------------------------
%	enumerate 及 item 的形狀
%----------------------------------------------------------------------------------------

\useinnertheme{rounded} % 圓球 (3D)
%\useinnertheme{circles} % 圓形 (2D)
%\useinnertheme{rectangles} % 方形
%\useinnertheme{inmargin} % 插入邊沿

%----------------------------------------------------------------------------------------
%	自訂 item 的顏色
%----------------------------------------------------------------------------------------

%\setbeamercolor{item projected}{bg=red}

%----------------------------------------------------------------------------------------
%	個人化的設置及細節調整
%----------------------------------------------------------------------------------------

% 設定頁面邊界
\setbeamersize{text margin left=0.8cm, text margin right=0.8cm}
\special{papersize=\the\paperwidth,\the\paperheight}
\providecommand{\tabularnewline}{\\}
}

%----------------------------------------------------------------------------------------
%	個人化的背景調控
%----------------------------------------------------------------------------------------

% 背景照片設置
%\setbeamertemplate{background}{\includegraphics[height=\paperheight]{Fig/Background.png}}

% 浮水印設定
\usebackgroundtemplate{%
	\tikz[overlay, remember picture] % 讓 logo 能每頁都顯示
	\node[opacity=0.3, below=-1.25cm, at=(current page.center)] % 調整透明度 (opacity) 及浮水印的位置
	{\includegraphics[scale = 0.14]{Fig/nthulogo.png}}; % 載入 logo 及調整大小
	}
